% theme: biology_and_cells
\MajorQuestion{生物と細胞}
\SetRowSpaceQanda

\Instruction{顕微鏡の各部と倍率に関する次の問いに答えなさい。}
\QQ{目を近づけてのぞくレンズを何というか。}{接眼レンズ}
\QQ{観察するものに近い側のレンズを何というか。}{対物レンズ}
\QQ{対物レンズを切り替えるために回す部分を何というか。}{レボルバー}
\QQ{プレパラートをのせる台を何というか。}{ステージ}
\QQ{顕微鏡の倍率は、何と何の倍率をかけて求めるか。}{接眼レンズと対物レンズ}

\Instruction{プレパラートのつくり方に関する次の問いに答えなさい。}
\QQ{観察するものをのせるガラス板を何というか。}{スライドガラス}
\QQ{観察するものの上にかぶせる薄いガラスを何というか。}{カバーガラス}
\QQ{観察するものに水を落とすときに使う器具を何というか。}{スポイト}
\QQ{カバーガラスをかぶせるときに使う器具を何というか。}{柄つき針}
\QQ{カバーガラスをかぶせるとき、入らないように注意するものは何か。}{気泡（空気の泡）}

\Instruction{細胞の基本的なつくりとはたらきに関する次の問いに答えなさい。}
\QQ{生物の体をつくる、小さな部屋のようなつくりを何というか。}{細胞}
\QQ{細胞の中にある、丸い形をした部分を何というか。}{核}
\QQ{細胞の中で、核のまわりにある部分を何というか。}{細胞質}
\QQ{細胞質の外側を包む薄い膜を何というか。}{細胞膜}
\QQ{細胞が酸素を取り入れ、二酸化炭素を出すはたらきを何というか。}{細胞の呼吸}

\Instruction{植物細胞と動物細胞、生物の体のつくりに関する次の問いに答えなさい。}
\QQ{植物細胞にあり、光合成を行う部分を何というか。}{葉緑体}
\QQ{植物細胞にあり、内部に液体を含む部分を何というか。}{液胞}
\QQ{植物細胞で、細胞膜の外側にある丈夫なつくりを何というか。}{細胞壁}
\QQ{体が1つの細胞だけでできている生物を何というか。}{単細胞生物}
\QQ{体が多くの細胞でできている生物を何というか。}{多細胞生物}
